\chapter*{前\hspace{1em}言}
\markboth{前\hspace{1em}言}{前\hspace{1em}言}
\phantomsection\label{ch:preface}

一本从零开始、一章只讲一件事的 Magel 入门书。

你不需要会 C，也不需要懂编译原理。这本书从头讲起，每章只解决一个问题，代码都能直接粘下来跑。读完它，你能独立写出结构清楚、资源管得干净的 Magel 程序，并且知道遇到问题时该去查哪一篇官方文档。

\begin{tip}{提示}
本书内容与 Magel 的当前\href{https://github.com/rabbitc999/Magel/tree/master/VERSION}{语言版本}和\href{https://github.com/rabbitc999/Magel/blob/master/doc/17-capabilities.md}{能力说明}保持同步。你可以按下方目录在线阅读，也可以下载 \href{https://github.com/rabbitc999/Magel/blob/master/doc/book/latex/book.pdf}{PDF 电子书}，离线阅读或打印。
\end{tip}

\section*{这本书写给谁}

\begin{itemize}

\item 刚开始学编程，想找一门“看得清楚”的语言入门；

\item 写过 C / Java / Go / Python，想快速知道 Magel 和它们哪里不一样；

\item 要维护一份逻辑、同时跑在服务端和嵌入式板子上，想知道这门语言能不能扛。

\end{itemize}

如果你只是想要一份语法清单，那你该看 \href{https://github.com/rabbitc999/Magel/blob/master/doc/04-language-reference.md}{语言参考}；这本书的目标不一样——它想让你读懂一门语言的设计，而不只是记住写法。

\section*{怎么读}

三条建议，按重要性排序：

\begin{enumerate}

\item 一定要动手。 每章的代码都短，敲一遍的收益远大于读十遍。随便建一个目录，把代码存成 \texttt{x.\allowbreak{}mgl}，然后 \texttt{magel run x.\allowbreak{}mgl}。

\item 按顺序读，别跳。 这本书是层层递进的：第 6 章解释为什么第 3 章的赋值那么“笨”，第 7 章的引用规则又是从第 6 章长出来的。跳着读，你会觉得语言在无理取闹。

\item 卡住了先看每章最后的「小结」和「下一步」。 那里写清了哪些内容本章刻意没讲、以及去下一章还是去查官方文档。

\end{enumerate}

书里的代码块分三类，看栅栏上的标注就知道：

\begin{center}
\begin{tabular}{ll}
\toprule
标注 & 含义 \\
\midrule
完整的 \texttt{main} 函数 & 可以整段复制出去跑，输出的内容就写在代码下面 \\
只有几行的片段 & 声明、表达式、语法示意，要放进程序里才能跑 \\
命令行 & 以 \texttt{magel} 开头，直接敲 \\
\bottomrule
\end{tabular}
\end{center}

\begin{tip}{提示}
Windows 上程序名是 \texttt{magel.\allowbreak{}exe}。在 Git Bash / MSYS 里写 \texttt{.\allowbreak{}/\allowbreak{}magel} 能自动找到它；光写 \texttt{magel} 会走 PATH 查找，可能找不到当前目录里这份。
\end{tip}

\section*{章节地图}

全书按第 1 至第 25 章顺序展开，难度逐步提升。第 6 至第 8 章集中介绍值语义、引用和资源管理，建议连贯阅读。

{\small\setlength{\tabcolsep}{4pt}
\begin{longtable}{>{\raggedright\arraybackslash}p{\dimexpr 0.229\linewidth-2\tabcolsep\relax}>{\raggedright\arraybackslash}p{\dimexpr 0.196\linewidth-2\tabcolsep\relax}>{\raggedright\arraybackslash}p{\dimexpr 0.555\linewidth-2\tabcolsep\relax}}
\toprule
章 & 标题 & 你会得到什么 \\
\midrule
\endhead
\hyperref[ch:01]{01} & 第一个程序 & 装好工具链，跑出 \texttt{Hello,\allowbreak{} Magel!}，知道 \texttt{run} 和 \texttt{check} 的区别 \\
\hyperref[ch:02]{02} & 猜数字 & 用一个小游戏把「变量、输入、循环、分支」串成一个能跑的程序 \\
\hyperref[ch:03]{03} & 变量与类型 & 变量、常量、\texttt{auto}，以及 Magel 为什么不肯帮你偷偷转类型 \\
\hyperref[ch:04]{04} & 函数 & 定义、参数、返回控制表达式、多值返回、重载与默认参数 \\
\hyperref[ch:05]{05} & 控制流 & \texttt{if} / \texttt{while} / \texttt{for} / \texttt{loop} / \texttt{for-in}，以及把 \texttt{loop} 当表达式用 \\
\hyperref[ch:06]{06} & 值语义 & 赋值就是复制。为什么不移动也能保证安全 \\
\hyperref[ch:07]{07} & 引用 & \texttt{mut\&} / \texttt{const\&}、没有 \texttt{null}、引用不能活得比目标久 \\
\hyperref[ch:08]{08} & 资源与作用域 & \texttt{Drop} / \texttt{Closeable} / \texttt{new} / \texttt{delete} / \texttt{move}，没有 GC 怎么活 \\
\hyperref[ch:09]{09} & 结构体 & 把字段和方法收进一个类型，构造器及 Equal/Hash/Clone 派生 \\
\hyperref[ch:10]{10} & 枚举与模式匹配 & 代数数据类型，以及让编译器帮你查漏的 \texttt{match} \\
\hyperref[ch:11]{11} & 容器 & \texttt{List} / \texttt{Map} / \texttt{Set}，以及容器也是值的后果 \\
\hyperref[ch:12]{12} & 字符串与文本 & 拼接、下标（按字节！）、切片，以及 \texttt{std:\allowbreak{}:\allowbreak{}text} \\
\hyperref[ch:13]{13} & 接口与运行期多态 & \texttt{interface}、默认实现、\texttt{std:\allowbreak{}:\allowbreak{}lang:\allowbreak{}:\allowbreak{}Any} \\
\hyperref[ch:14]{14} & 泛型 & 泛型函数与结构体、约束、编译期单态化 \\
\hyperref[ch:15]{15} & 函数值与闭包 & \texttt{fn} 类型、具名局部闭包与四种捕获方式 \\
\hyperref[ch:16]{16} & 错误处理 & \texttt{Option} / \texttt{Result} / \texttt{?}，以及 \texttt{assert} 该用在哪 \\
\hyperref[ch:17]{17} & 迭代器 & \texttt{std:\allowbreak{}:\allowbreak{}iter} 的惰性链：\texttt{range} / \texttt{map} / \texttt{filter} / \texttt{collect} \\
\hyperref[ch:18]{18} & 模块与工程 & \texttt{module} / \texttt{import}、\texttt{magel new}、\texttt{magel.\allowbreak{}toml}、测试与格式化 \\
\hyperref[ch:19]{19} & 预处理与宏 & 条件编译、文件包含与有序多括号文本宏 \\
\hyperref[ch:20]{20} & 注解与反射 & \texttt{@annotation}、运行期查元数据、按名调用 \\
\hyperref[ch:21]{21} & 与 C 互操作 & \texttt{abi} 块，调 C 的库、把自己的函数交给 C \\
\hyperref[ch:22]{22} & 交付与部署 & \texttt{.\allowbreak{}svmb}、AOT、符号压缩、嵌入式与看门狗 \\
\hyperref[ch:23]{23} & 综合项目：命令行待办清单 & 把前面二十章的东西拼成一个能用的程序 \\
\hyperref[ch:24]{24} & 学生信息管理系统 & MySQL、HTTP API、跨源前端、会话与权限，配真实界面截图 \\
\hyperref[ch:25]{25} & 单片机上的 Magel 解释器 & H743 共享 VM、字节码部署、内存预算、外设胶水与退出恢复 \\
\bottomrule
\end{longtable}
}

\subsection*{附录}

\begin{center}
\begin{tabular}{ll}
\toprule
附录 & 内容 \\
\midrule
\hyperref[ch:appA]{附录 A} & 常见报错、诊断码与处理办法 \\
\hyperref[ch:appB]{附录 B} & 按能力域列出的标准库模块，以及去哪儿查细节 \\
\bottomrule
\end{tabular}
\end{center}

\section*{和官方文档的关系}

这本书不替代官方文档，它是官方文档的导读。

\begin{center}
\begin{tabular}{ll}
\toprule
你想 & 去哪 \\
\midrule
跟着书学语言 & 就是这儿 \\
查某个语法怎么写 & \href{https://github.com/rabbitc999/Magel/blob/master/doc/04-language-reference.md}{语言参考} \\
理解这些设计为什么这么做 & \href{https://github.com/rabbitc999/Magel/blob/master/doc/03-core-concepts.md}{核心概念} \\
5 分钟跑起来 & \href{https://github.com/rabbitc999/Magel/blob/master/doc/01-quickstart.md}{快速开始} \\
查某个标准库函数 & \href{https://github.com/rabbitc999/Magel/blob/master/doc/library/README.md}{库函数 API 参考} \\
看官方给的可运行样例 & \href{https://github.com/rabbitc999/Magel/tree/master/example/samples}{\texttt{example/\allowbreak{}samples/\allowbreak{}}} \\
\bottomrule
\end{tabular}
\end{center}

本书正文里凡是引用官方文档的地方，都给了相对链接，可以随时跳过去。

\section*{开始之前}

\begin{itemize}

\item 工具链：见 \hyperref[ch:01]{第 1 章}。

\item 编辑器：推荐装 \href{https://github.com/rabbitc999/Magel/tree/master/vscode}{\texttt{vscode/\allowbreak{}}} 下的 Magel 插件，有语法高亮、补全、跳转与断点调试；不装也能读完整本书。

\item 心里准备：第 6 到第 8 章会请你改掉一些习惯。别急着问“为什么不直接给我”，作者在那三章里把取舍写清楚了。

\end{itemize}

准备好了就翻到 \hyperref[ch:01]{第 1 章}。
